% Part: first-order-logic
% Chapter: natural-deduction
% Section: soundness

\documentclass[../../../include/open-logic-section]{subfiles}

\begin{document}

\iftag{FOL}
      {\olfileid{fol}{ntd}{sou}}
      {\olfileid{pl}{ntd}{sou}}

\olsection{Kasahihan}

\begin{explain}
Sistem !!^a{derivation}, kayata deduksi natural, iku \emph{sah} yen
ora bisa !!{derive} bab kang sejatine ora ndherek. Mula kasahihan iku
sawijining sipat aman kang dijamin kanggo sistem
!!{derivation}. Gumantung marang sipat teoretis-bukti kang dirembug,
kita kepengin ngerti, upamane, yen
\begin{enumerate}
\item saben !!{derivable} !!{sentence} iku
  \iftag{FOL}{valid}{tautologi};
\item yen !!a{sentence} bisa !!{derivable} saka ukara-ukara liyane,
  ukara iku uga dadi akibate;
\item yen sawijining himpunan !!{sentence}s ora konsisten, himpunan
  iku ora bisa dicukupi.
\end{enumerate}
Iki sipat-sipat wigati saka sistem !!a{derivation}. Yen salah sijine
ora lumaku, sistem !!{derivation} kasebut cacad---sistem iku bakal
!!{derive} kakehan. Mula, netepake kasahihan sistem !!{derivation}
iku wigati banget.
\end{explain}

\begin{thm}[Kasahihan]
\ollabel{thm:soundness}
Yen $!A$ bisa !!{derivable} saka asumsi !!{undischarged} $\Gamma$,
mula $\Gamma \Entails !A$.
\end{thm}

\begin{proof}
Ayo $\delta$ dadi !!a{derivation} saka $!A$. Kita lumaku kanthi
induksi ing cacah inferensi ing~$\delta$.

Kanggo basis induksi, kita nuduhake pratelan iki nalika cacah
inferensi yaiku~$0$. Ing kasus iki, $\delta$ mung dumadi saka siji
!!{sentence}~$!A$, yaiku sawijining asumsi. Asumsi kasebut
!!{undischarged}, amarga asumsi mung bisa !!{discharged} dening
inferensi, nanging ora ana inferensi. Dadi, saben
\iftag{FOL}{!!{structure}~$\Struct{M}$}{!!{valuation}~$\pAssign{v}$}
kang nyukupi kabeh asumsi !!{undischarged} ing bukti kasebut uga
nyukupi~$!A$.

Saiki kanggo langkah induktif. Upamana $\delta$ ngemot~$n$ inferensi.
Premis siji utawa luwih saka inferensi paling ngisor iku !!{derive}d
nganggo sub-!!{derivation}s kang saben-sabene ngemot inferensi kurang
saka~$n$. Kita nganggep hipotesis induksi: premis-premis saka
inferensi paling ngisor ndherek saka asumsi !!{undischarged} ing
sub-!!{derivation}s kang mungkasi ing premis kasebut. Kita kudu
nuduhake yen kesimpulan~$!A$ ndherek saka asumsi !!{undischarged} ing
kabeh bukti.

Kita mbedakake kasus adhedhasar jinis inferensi paling ngisor. Sepisan,
kita nimbang inferensi-inferensi kang mung duwe siji premis.

\begin{enumerate}
\item Upamana inferensi pungkasan iku \Intro{\lnot}:
  !!{derivation} kasebut awujud
  \begin{prooftree}
    \AxiomC{$\Gamma, \Discharge{!A}{n}$}
    \RightLabel{$\delta_1$}
    \DeduceC{$\lfalse$}
    \DischargeRule{\Intro{\lnot}}{n}
    \UnaryInfC{$\lnot !A$}
  \end{prooftree}
  Miturut hipotesis induksi, $\lfalse$ ndherek saka asumsi
  !!{undischarged} $\Gamma \cup \{!A\}$ ing~$\delta_1$. Nimbang
  \iftag{FOL}{!!a{structure}~$\Struct{M}$}{!!a{valuation}~$\pAssign{v}$}.
  Kita kudu nuduhake yen, yen
  $\iftag{FOL}{\Sat{M}}{\pSat{v}}{\Gamma}$, mula
  $\iftag{FOL}{\Sat{M}}{\pSat{v}}{\lnot !A}$. Kanggo reductio,
  upamana $\iftag{FOL}{\Sat{M}}{\pSat{v}}{\Gamma}$, nanging
  $\iftag{FOL}{\Sat/{M}}{\pSat/{v}}{\lnot !A}$, yaiku
  $\iftag{FOL}{\Sat{M}}{\pSat{v}}{!A}$. Iki bakal ateges yen
  $\iftag{FOL}{\Sat{M}}{\pSat{v}}{\Gamma \cup \{!A\}}$. Iki
  mbantah hipotesis induksi kita. Dadi,
  $\iftag{FOL}{\Sat{M}}{\pSat{v}}{\lnot !A}$.

\item Inferensi pungkasan iku \Elim{\land}: Ana rong varian: $!A$
  utawa $!B$ bisa diturunake saka premis $!A \land !B$. Nimbang kasus
  kapisan. !!{derivation}~$\delta$ katon kaya mangkene:
  \begin{prooftree}
    \AxiomC{$\Gamma$}
    \RightLabel{$\delta_1$}
    \DeduceC{$!A \land !B$}
    \RightLabel{\Elim{\land}}
    \UnaryInfC{$!A$}
  \end{prooftree}
  Miturut hipotesis induksi, $!A \land !B$ ndherek saka asumsi
  !!{undischarged}~$\Gamma$ ing~$\delta_1$. Nimbang
  !!a{structure}~\iftag{FOL}{$\Struct{M}$}{$\pAssign{v}$}. Kita kudu
  nuduhake yen, yen $\iftag{FOL}{\Sat{M}}{\pSat{v}}{\Gamma}$, mula
  $\iftag{FOL}{\Sat{M}}{\pSat{v}}{!A}$. Upamana
  $\iftag{FOL}{\Sat{M}}{\pSat{v}}{\Gamma}$. Miturut hipotesis induksi
  kita ($\Gamma \Entails !A \land !B$), kita ngerti yen
  $\iftag{FOL}{\Sat{M}}{\pSat{v}}{!A \land !B}$. Miturut definisi,
  $\iftag{FOL}{\Sat{M}}{\pSat{v}}{!A \land !B}$ yen lan mung yen
  $\iftag{FOL}{\Sat{M}}{\pSat{v}}{!A}$ lan
  $\iftag{FOL}{\Sat{M}}{\pSat{v}}{!B}$. (Kasus nalika $!B$
  diturunake saka $!A \land !B$ ditangani kanthi cara kang padha.)

\item Inferensi pungkasan iku \Intro{\lor}: Ana rong varian: $!A
  \lor !B$ bisa diturunake saka premis~$!A$ utawa premis~$!B$.
  Nimbang kasus kapisan. !!{derivation} kasebut awujud
  \begin{prooftree}
    \AxiomC{$\Gamma$}
    \RightLabel{$\delta_1$}
    \DeduceC{$!A$}
    \RightLabel{\Intro{\lor}}
    \UnaryInfC{$!A \lor !B$}
  \end{prooftree}
  Miturut hipotesis induksi, $!A$ ndherek saka asumsi
  !!{undischarged}~$\Gamma$ ing~$\delta_1$. Nimbang
  \iftag{FOL}{!!a{structure}~$\Struct{M}$}{!!a{valuation}~$\pAssign{v}$}.
  Kita kudu nuduhake yen, yen
  $\iftag{FOL}{\Sat{M}}{\pSat{v}}{\Gamma}$, mula
  $\iftag{FOL}{\Sat{M}}{\pSat{v}}{!A \lor !B}$. Upamana
  $\iftag{FOL}{\Sat{M}}{\pSat{v}}{\Gamma}$; banjur
  $\iftag{FOL}{\Sat{M}}{\pSat{v}}{!A}$ amarga $\Gamma \Entails !A$
  (hipotesis induksi). Mula uga mesthi
  $\iftag{FOL}{\Sat{M}}{\pSat{v}}{!A \lor !B}$. (Kasus nalika $!A
  \lor !B$ diturunake saka~$!B$ ditangani kanthi cara kang padha.)

\item Inferensi pungkasan iku \Intro{\lif}: $!A \lif !B$
  diturunake saka subbukti kang duwe asumsi~$!A$ lan kesimpulan~$!B$,
  yaiku
  \begin{prooftree}
    \AxiomC{$\Gamma, \Discharge{!A}{n}$}
    \RightLabel{$\delta_1$}
    \DeduceC{$!B$}
    \DischargeRule{\Intro{\lif}}{n}
    \UnaryInfC{$!A \lif !B$}
  \end{prooftree}
  Miturut hipotesis induksi, $!B$ ndherek saka asumsi
  !!{undischarged} ing~$\delta_1$, yaiku $\Gamma \cup \{!A\}
  \Entails !B$. Nimbang
  \iftag{FOL}{!!a{structure}~$\Struct{M}$}{!!a{valuation}~$\pAssign{v}$}.
  Asumsi !!{undischarged} ing~$\delta$ mung $\Gamma$, amarga $!A$ wis
  dibubarake ing inferensi pungkasan. Dadi kita kudu nuduhake yen
  $\Gamma \Entails !A \lif !B$. Kanggo reductio, upamana kanggo
  sawijining
  \iftag{FOL}{!!{structure}~$\Struct{M}$}{!!{valuation}~$\pAssign{v}$},
  $\iftag{FOL}{\Sat{M}}{\pSat{v}}{\Gamma}$ nanging
  $\iftag{FOL}{\Sat/{M}}{\pSat/{v}}{!A \lif !B}$. Dadi,
  $\iftag{FOL}{\Sat{M}}{\pSat{v}}{!A}$ lan
  $\iftag{FOL}{\Sat/{M}}{\pSat/{v}}{!B}$. Nanging miturut hipotesis,
  $!B$ iku akibat saka $\Gamma \cup \{!A\}$, yaiku
  $\iftag{FOL}{\Sat{M}}{\pSat{v}}{!B}$, kang dadi kontradiksi. Mula,
  $\Gamma \Entails !A \lif !B$.

\item Inferensi pungkasan iku \FalseInt: Ing kene, $\delta$
  mungkasi nganggo
  \begin{prooftree}
    \AxiomC{$\Gamma$}
    \RightLabel{$\delta_1$}
    \DeduceC{$\lfalse$}
    \RightLabel{\FalseInt}
    \UnaryInfC{$!A$}
  \end{prooftree}
  Miturut hipotesis induksi, $\Gamma \Entails \lfalse$. Kita kudu
  nuduhake yen $\Gamma \Entails !A$. Upamana ora; banjur kanggo
  sawijining~$\iftag{FOL}{\Struct{M}}{\pAssign{v}}$ kita duwe
    $\iftag{FOL}{\Sat{M}}{\pSat{v}}{\Gamma}$ lan
    $\iftag{FOL}{\Sat/{M}}{\pSat/{v}}{!A}$. Nanging kita tansah duwe
    $\iftag{FOL}{\Sat/{M}}{\pSat/{v}}{\lfalse}$, mula iki bakal ateges
    yen $\Gamma \Entails/ \lfalse$, kang mbantah hipotesis induksi.

\item Inferensi pungkasan iku \FalseCl: Gladhen.

\iftag{FOL}{
\item Inferensi pungkasan iku \Intro{\lforall}: Banjur $\delta$
  awujud
  \begin{prooftree}
    \AxiomC{$\Gamma$}
    \RightLabel{$\delta_1$}
    \DeduceC{$!A(a)$}
    \RightLabel{\Intro{\lforall}}
    \UnaryInfC{$\lforall[x][!A(x)]$}
  \end{prooftree}
  Miturut hipotesis induksi, premis $!A(a)$ iku akibat saka asumsi
  !!{undischarged} $\Gamma$. Nimbang sawijining struktur,
  $\Struct{M}$, kang $\Sat{M}{\Gamma}$. Kita kudu nuduhake yen
  $\Sat{M}{\lforall[x][!A(x)]}$. Amarga $\lforall[x][!A(x)]$ iku
  !!a{sentence}, iki ateges kita kudu nuduhake yen kanggo saben
  penetapan variabel~$s$, $\Sat{M}{!A(x)}[s]$
  (\olref[syn][ass]{prop:sat-quant}). Amarga $\Gamma$ mung dumadi saka
  ukara, $\Sat{M}{!B}[s]$ kanggo kabeh $!B \in \Gamma$ miturut
  \olref[syn][sat]{defn:satisfaction}. Ayo $\Struct{M'}$ padha karo
  $\Struct{M}$ kajaba $\Assign{a}{M'} = s(x)$. Amarga $a$ ora muncul
  ing~$\Gamma$, $\Sat{M'}{\Gamma}$ miturut
  \olref[syn][ext]{cor:extensionality-sent}. Amarga $\Gamma \Entails
  !A(a)$, $\Sat{M'}{!A(a)}$. Amarga $!A(a)$ iku !!a{sentence},
  $\Sat{M'}{!A(a)}[s]$ miturut
  \olref[syn][ass]{prop:sentence-sat-true}. $\Sat{M'}{!A(x)}[s]$ yen
  lan mung yen $\Sat{M'}{!A(a)}$ miturut
  \olref[syn][ext]{prop:ext-formulas} (elinga yen $!A(a)$ mung
  $\Subst{!A(x)}{a}{x}$). Dadi, $\Sat{M'}{!A(x)}[s]$. Amarga $a$ ora
  muncul ing~$!A(x)$, miturut
  \olref[syn][ext]{prop:extensionality}, $\Sat{M}{!A(x)}[s]$. Nanging
  $s$ iku penetapan variabel sawenang-wenang, mula
  $\Sat{M}{\lforall[x][!A(x)]}$.

\item Inferensi pungkasan iku \Intro{\lexists}: Gladhen.

\item Inferensi pungkasan iku \Elim{\forall}: Gladhen.
}{}
\end{enumerate}

Saiki ayo nimbang inferensi-inferensi kang bisa duwe sawetara premis:
\Elim{\lor}, \Intro{\land},
\iftag{FOL}{\Elim{\lif}, lan \Elim{\lexists}}{lan \Elim{\lif}}.
\begin{enumerate}

\item Inferensi pungkasan iku \Intro{\land}. $!A \land !B$
  diturunake saka premis $!A$ lan $!B$, lan $\delta$ awujud
  \begin{prooftree}
    \AxiomC{$\Gamma_1$}
    \RightLabel{$\delta_1$}
    \DeduceC{$!A$}
    \AxiomC{$\Gamma_2$}
    \RightLabel{$\delta_2$}
    \DeduceC{$!B$}
    \RightLabel{\Intro{\land}}
    \BinaryInfC{$!A \land !B$}
  \end{prooftree}
  Miturut hipotesis induksi, $!A$ ndherek saka asumsi
  !!{undischarged}~$\Gamma_1$ ing~$\delta_1$, lan $!B$ ndherek saka
  asumsi !!{undischarged}~$\Gamma_2$ ing~$\delta_2$. Asumsi
  !!{undischarged} ing~$\delta$ yaiku $\Gamma_1 \cup \Gamma_2$, mula
  kita kudu nuduhake yen $\Gamma_1 \cup \Gamma_2 \Entails !A \land
  !B$. Nimbang
  \iftag{FOL}{!!a{structure}~$\Struct{M}$}{!!a{valuation}~$\pAssign{v}$}
  kanthi $\iftag{FOL}{\Sat{M}}{\pSat{v}}{\Gamma_1 \cup \Gamma_2}$.
  Amarga $\iftag{FOL}{\Sat{M}}{\pSat{v}}{\Gamma_1}$, mesthi
  $\iftag{FOL}{\Sat{M}}{\pSat{v}}{!A}$ awit $\Gamma_1 \Entails !A$,
  lan amarga $\iftag{FOL}{\Sat{M}}{\pSat{v}}{\Gamma_2}$,
  $\iftag{FOL}{\Sat{M}}{\pSat{v}}{!B}$ awit $\Gamma_2 \Entails !B$.
  Kalorone bebarengan menehi
  $\iftag{FOL}{\Sat{M}}{\pSat{v}}{!A \land !B}$.

\item Inferensi pungkasan iku \Elim{\lor}: Gladhen.

\item Inferensi pungkasan iku \Elim{\lif}. $!B$ diturunake saka
  premis $!A \lif !B$ lan~$!A$. !!{derivation}~$\delta$ katon kaya
  mangkene:
  \begin{prooftree}
    \AxiomC{$\Gamma_1$}
    \RightLabel{$\delta_1$}
    \DeduceC{$!A \lif !B$}
    \AxiomC{$\Gamma_2$}
    \RightLabel{$\delta_2$}
    \DeduceC{$!A$}
    \RightLabel{\Elim{\lif}}
    \BinaryInfC{$!B$}
  \end{prooftree}
  Miturut hipotesis induksi, $!A \lif !B$ ndherek saka asumsi
  !!{undischarged}~$\Gamma_1$ ing~$\delta_1$, lan $!A$ ndherek saka
  asumsi !!{undischarged}~$\Gamma_2$ ing~$\delta_2$. Nimbang
  \iftag{FOL}{!!a{structure}~$\Struct{M}$}{!!a{valuation}~$\pAssign{v}$}.
  Kita kudu nuduhake yen, yen
  $\iftag{FOL}{\Sat{M}}{\pSat{v}}{\Gamma_1 \cup \Gamma_2}$, mula
  $\iftag{FOL}{\Sat{M}}{\pSat{v}}{!B}$. Upamana
  $\iftag{FOL}{\Sat{M}}{\pSat{v}}{\Gamma_1 \cup \Gamma_2}$. Amarga
  $\Gamma_1 \Entails !A \lif !B$,
  $\iftag{FOL}{\Sat{M}}{\pSat{v}}{!A \lif !B}$. Amarga $\Gamma_2
  \Entails !A$, kita duwe $\iftag{FOL}{\Sat{M}}{\pSat{v}}{!A}$.
  Iki ateges $\iftag{FOL}{\Sat{M}}{\pSat{v}}{!B}$ (amarga yen
  $\iftag{FOL}{\Sat/{M}}{\pSat/{v}}{!B}$, awit
  $\iftag{FOL}{\Sat{M}}{\pSat{v}}{!A}$, kita bakal duwe
  $\iftag{FOL}{\Sat/{M}}{\pSat/{v}}{!A \lif !B}$, kang mbantah
  $\iftag{FOL}{\Sat{M}}{\pSat{v}}{!A \lif !B}$).

\item Inferensi pungkasan iku \Elim{\lnot}: Gladhen.

\tagitem{FOL}{Inferensi pungkasan iku \Elim{\lexists}: Gladhen.}{}
\end{enumerate}
\end{proof}

\tagprob{FOL}
\begin{prob}
Jangkepana bukti \olref[fol][ntd][sou]{thm:soundness}.
\end{prob}
\tagendprob

\tagprob{notFOL}
\begin{prob}
Jangkepana bukti \olref[pl][ntd][sou]{thm:soundness}.
\end{prob}
\tagendprob

\begin{cor}
\ollabel{cor:weak-soundness}
Yen $\Proves !A$, mula $!A$ iku \iftag{FOL}{valid}{tautologi}.
\end{cor}

\begin{cor}
\ollabel{cor:consistency-soundness}
Yen $\Gamma$ bisa dicukupi, mula himpunan iku konsisten.
\end{cor}

\begin{proof}
Kita mbuktekake kontraposisine. Upamana $\Gamma$ ora konsisten.
Banjur $\Gamma \Proves \lfalse$, yaiku ana !!a{derivation} saka
$\lfalse$ adhedhasar asumsi !!{undischarged} ing~$\Gamma$. Miturut
\olref{thm:soundness}, saben
\iftag{FOL}{!!{structure}~$\Struct{M}$}{!!{valuation}~$\pAssign{v}$}
kang nyukupi $\Gamma$ mesthi nyukupi $\lfalse$. Amarga
$\iftag{FOL}{\Sat/{M}}{\pSat/{v}}{\lfalse}$ kanggo saben
\iftag{FOL}{!!{structure}~$\Struct{M}$}{!!{valuation}~$\pAssign{v}$},
ora ana \iftag{FOL}{$\Struct{M}$}{$\pAssign{v}$} kang bisa nyukupi
$\Gamma$, yaiku $\Gamma$ ora bisa dicukupi.
\end{proof}

\end{document}
